\chapter{简历表达与面试复盘}

\begin{beginnerbox}
简历不是把仓库里的技术名词全部塞进去，而是用少量可证明的句子，让面试官知道：你为什么做、解决谁的问题、你负责什么、系统怎样工作、证据是什么、边界在哪里。只有自己能解释、复现和修改的内容，才适合写成“我设计”或“我实现”。

学习时先不要背整段答案。先用自己的话说清楚每个阶段解决的业务问题，再把本章的回答模板当作整理表达的工具。
\end{beginnerbox}

\section{先确定项目的一句话定位}

推荐项目名：

\begin{quote}
ECom Review Intelligence Agent：面向品牌方产品/类目运营的评论证据决策原型
\end{quote}

一句话介绍：

\begin{quote}
我围绕 SKU 诊断、竞品比较和问题分诊，构建了一个由本地证据工具掌握事实与执行权、DeepSeek 负责受约束规划和表达、Verifier 核对业务字段与引用的 LangGraph Agent。
\end{quote}

这句话比“做了一个电商 RAG 聊天机器人”更准确，因为项目包含数据合同、商品授权、确定性工具、澄清恢复、逐字段验证、评测和安全发布，而不是只有检索与生成。

\section{为什么做这个项目}

面试时可以分三层回答：

\begin{enumerate}
  \item \textbf{个人背景：}我在运营中心实习，也关注广州电商和 AI 应用岗位，希望项目能同时体现业务理解与工程能力；
  \item \textbf{业务问题：}运营阅读大量评论后，仍需要把问题信号整理成能回到证据的单品诊断、竞品比较和跨团队分诊；
  \item \textbf{技术问题：}如果只让模型自由总结，商品范围、引用、工具权限和失败状态都难以审计，因此我把事实、执行权和校验留在本地代码。
\end{enumerate}

\begin{boundarybox}
广州和运营实习是选题动机，不是完整市场证明。可以说“这个方向与我的目标岗位和经历匹配”，不能说“广州所有电商公司都需要”或“市场已经验证付费意愿”。
\end{boundarybox}

\section{通用简历版本}

可根据版面选三到四条：

\begin{itemize}
  \item 面向品牌方产品/类目运营，定义 SKU 诊断、竞品比较和问题分诊三条评论证据工作流，将效率、决策支持与治理价值假设和已验证工程结果分开。
  \item 从冻结数据快照、VOC 证据合同到 SQLite/FTS5 本地证据库建立可重算 lineage，并用稳定本地引用、精确商品关联和公开边界支持可复核分析。
  \item 构建受约束 LangGraph 单 Agent，结合严格 Structured Output、本地确定性工具、一次澄清恢复、checkpoint、请求预算、安全证据卡、逐字段 Verifier 和失败关闭。
  \item 完成 30 个冻结离线 Agent 案例及真实 DeepSeek 15 乘 2 有界评估，并将 Agent 封装为 loopback-only typed API 与桌面/移动运营工作台；P7 真人运营试用仍待完成。
\end{itemize}

\section{按岗位调整重点}

\subsection{AI 产品或数据产品岗位}

强调：为什么选这个问题、主用户是谁、三工作流如何从需求到指标、为什么控制 MVP、怎样区分价值假设与工程事实、下一步怎样做用户验证。

可替换 bullet：

\begin{quote}
设计“原流程与工作台同任务比较”的 P7 用户试用方案，计划记录首次有效证据、引用复核、采纳/修改/拒绝和失败原因；在目标用户试用前不宣称节省工时或改善决策。
\end{quote}

\subsection{AI 应用或 Agent 岗位}

强调：严格 DTO、tool/authority 分离、LangGraph 状态、interrupt/resume、checkpoint、Verifier、fake/replay/live、离线与真实模型 gate、公开/私有 Trace 和失败原子发布。

可替换 bullet：

\begin{quote}
修复真实模型“JSON 合法但业务字段未授权”的问题，通过动态收窄 schema 与本地逐字段 Verifier 核对 workflow、别名、极性、强度和引用，失败只修复一次后确定性降级。
\end{quote}

\subsection{电商运营或 VOC 岗位}

强调：评论证据怎样支持单品体验问题整理、竞品同维度比较、质量/客服协同和下一步核查；避免大段讲框架内部实现。

可替换 bullet：

\begin{quote}
围绕噪声、可靠性、清洁维护等八类 VOC 维度，将评论问题信号整理为可追溯证据卡和进一步核查建议，明确不把评论频次外推为缺陷率、销量损失或因果结论。
\end{quote}

\section{30 秒介绍}

\begin{quote}
我做的是一个面向品牌方产品和类目运营的评论证据 Agent，支持 SKU 诊断、竞品比较和问题分诊。它不是让模型直接掌握事实，而是由本地工具取证，模型做受约束规划和表达，Verifier 再核对引用和业务规则。P0 到 P6 已通过，P7 的 synthetic 桌面/移动验证和 P6-backed 本地三工作流回归也已通过；但目标运营用户尚未试用，所以我不宣称已经产生业务效率或降本增效。
\end{quote}

\section{1 分钟介绍}

\begin{quote}
这个项目来自我的运营中心实习和广州电商求职方向。我发现普通评论摘要很难回答“结论能否回到证据”和“模型有没有越权”，所以我把项目定义成 B 端内部决策支持原型，而不是消费者聊天助手。

数据层先冻结来源、品类范围和八类 VOC 合同，再用 SQLite/FTS5 建本地证据库。Agent 层用 LangGraph 管理诊断、比较、分流和一次澄清恢复；外部 DeepSeek 只看到匿名安全卡，本地 resolver、router、tool ledger 和 Verifier 分别掌握商品、工具、次数和事实授权。

验证上，我保留了检索负结果和 Bad Case，完成离线 30 案例与真实模型 15 乘 2 有界 gate。P7 工作台工程路径已验证，但真人运营试用仍未完成，所以当前结论是“工程闭环可复核”，不是“业务效果已证明”。
\end{quote}

\section{3 分钟介绍}

\begin{quote}
我希望这个项目同时能回答业务和技术两个问题。业务上，品牌方产品或类目运营需要从大量评论中整理单品问题、竞品差异和跨团队核查优先级；技术上，直接让大模型总结会把商品范围、工具权限、引用和失败状态都交给一个概率模型，难以审计。

因此我先没有做 Agent。P0 和 P1 验证数据源并冻结快照，P2 通过流式审计选择 Ice Makers 作为两个月 MVP，P3 冻结 VOC、极性、强度和引用合同，P4 建 SQLite/FTS5 确定性证据库，P5 用冻结 silver 案例比较检索基线并保留 BM25 没有独立胜出的负结果。这样 P6 的模型只能消费已经定义的证据，而不能临时改变口径。

P6 使用 LangGraph 管理状态和分支。输入先经 Pydantic，模型提交受约束计划，本地 resolver 决定商品，router 只授权一个业务工具，ledger 限制执行次数。工具返回后，本地代码把结果变成匿名安全卡，DeepSeek 提交结构化草稿，Verifier 再逐字段核对 workflow、商品别名、aspect、卡片、极性、强度和引用。标题歧义时 interrupt 并通过 UUID4 token resume；模型失败、证据不足或二次拒绝时确定性降级。

我遇到过两个关键问题：一是离线评测曾可能零引用真空通过，后来改成 expected、tool、rendered 引用三方相等并逐一回库；二是真实 DeepSeek 即使返回合法 JSON，也会改写业务字段，因此增加动态 schema 与本地 Verifier。最后完成离线 30 案例、真实模型 15 个案例各两次的有界 gate 和完成审计。

P7 把这张图封装成 loopback-only typed API 和运营工作台，synthetic 浏览器与本地三工作流接线已经验证。现在最重要的下一步不是再加框架，而是让目标运营用户完成同任务的手工基线和工作台试用，记录首次有效证据、引用复核、采纳和失败原因。只有拿到这些证据后，才能讨论是否真正节省时间或改善决策。
\end{quote}

\section{高频业务问题}

\subsection{Q1：项目的业务价值到底是什么？}

\textbf{回答框架：}当前已实现的是把评论整理为可追溯的诊断、比较和分诊证据，可能减少重复筛选并改善跨团队沟通；工程治理价值已经验证，效率和决策改善仍待目标用户试用。不要直接说“降本增效”。

\subsection{Q2：为什么主要面向 B 端，不直接做消费者助手？}

\textbf{回答框架：}我的经历和目标岗位更接近运营；B 端任务需要证据、协作和审计，能体现产品与工程闭环。底层能力未来可支持 C 端，但当前 MVP 不做推荐、价格和库存决策。

\subsection{Q3：运营为什么不用 Excel 或 BI？}

\textbf{回答框架：}Excel/BI 很适合结构化指标，本项目补的是评论文本到证据卡、受控工作流和可复核自然语言之间的环节。是否值得替代现有流程仍需要用户试验，不能假设 Agent 一定更好。

\subsection{Q4：怎样验证业务价值？}

\textbf{回答框架：}对同一任务先做原流程，再做工作台流程，记录完成时间、首次有效证据、引用复核、采纳/修改/拒绝和失败原因。单个用户只能形成定性观察，扩大结论需要更多角色和任务。

\section{高频技术问题}

\subsection{Q5：为什么这不是普通 RAG？}

检索只是其中一层。系统还包含商品授权、三工作流、澄清恢复、工具 ledger、逐字段 Verifier、请求预算、双轨 Trace、离线/live gate 和失败原子发布。

\subsection{Q6：为什么用 LangGraph？不用行不行？}

固定脚本也能实现简单主路径。选择 LangGraph 是因为项目有澄清暂停、恢复、一次修复、确定性降级和 checkpoint，需要把节点和状态显式化。LangGraph 负责编排，不替代业务校验或安全策略。

\subsection{Q7：Structured Output 已经有 schema，为什么还要 Verifier？}

schema 只能保证字段与类型可以解析，不能保证这些值来自本次授权证据。真实模型曾返回结构合法但极性或强度未授权的结果，所以本地 Verifier 必须逐字段比较。

\subsection{Q8：怎样防止幻觉？}

不承诺消灭幻觉。调用前只发送匿名安全卡；模型无数据库和任意工具权限；生成后逐字段校验，引用三方相等并回库；失败只修复一次，再失败则安全降级。

\subsection{Q9：为什么不用向量数据库？}

P5 的当前冻结案例没有证明复杂检索的独立收益，而且 silver 集有词法循环偏向。应先建立人工 relevance 数据集，再决定是否引入向量检索，而不是为简历堆技术。

\subsection{Q10：Fake、Replay 和 Live 有什么区别？}

fake 用于程序内故障注入，replay 用于冻结复现，live 才验证真实供应商协议和行为。任何一种模式通过，都不能替代另外两种想证明的内容。

\subsection{Q11：15 乘 2 能证明生产可靠性吗？}

不能。它证明当前代码、配置、冻结案例和当前供应商版本下的有界行为，不覆盖开放输入、并发、长期网络波动、未来模型版本和生产 SLA。

\subsection{Q12：checkpoint 里存了什么？}

存可序列化的 Agent 状态和预算，不存 live client 或不可序列化对象。恢复时重新注入运行时依赖，token 必须属于原线程且预算只能单调增加。

\section{失败案例追问}

\subsection{Q13：最有价值的失败是什么？}

可以回答离线 false PASS：早期引用可能在空集合上真空通过，安全计数也可能由场景推导。我先写反例测试，再要求 expected/tool/rendered 三方相等、回真实库验证，并让每个业务与安全预期成为阻断门。这个案例体现了我对“测试通过能证明什么”的重新理解。

\subsection{Q14：真实 API 接入最难的地方是什么？}

不是把 HTTP 请求发出去，而是区分协议格式和业务授权。JSON mode 解决可解析问题，强制 tool call 解决草稿结构，动态 schema 和本地 Verifier 才解决当前证据授权。

\subsection{Q15：P7 为什么还没有完成？}

工程自动化、桌面移动视口和 P6-backed 本地接线已经验证，但验收定义还要求目标运营用户的结构化试用。开发者自动化不能替代用户是否理解、采纳和完成任务的证据。

\section{压力问题和诚实回答}

\subsection{“你是不是用 AI 写的？”}

如实说明 AI 用于生成、审查、测试建议和文档整理。个人贡献用四件事界定：能否解释设计取舍、能否从干净环境复现、能否在反例下修改、能否对结论边界负责。不要把 AI 生成的每行代码都说成纯手写。

\subsection{“这个项目真的有人用吗？”}

当前没有目标运营用户完成正式结构化试用。可以展示 synthetic 和 local 工程证据，但不能把自己或自动化浏览器称为真实用户。下一步试用计划已经定义。

\subsection{“业务指标是多少？”}

当前没有可靠的工时节省、采纳率或决策质量指标。开发者时延和测试通过数不是业务指标。若面试官需要，可以说明准备怎样测，而不是现场编造数字。

\subsection{“你不会的问题怎么办？”}

先说明目前能确认的证据，再指出缺少什么信息。例如：“我目前没有多租户和生产并发证据，所以不能判断部署规模；我会先定义流量、隔离和 SLA，再做容量与故障测试。”证据不足不等于结论错误，但不能装成已确认。

\section{当前系统边界清单}

面试中应主动说明当前不是成熟企业 SaaS，仍缺：

\begin{itemize}
  \item 多租户、RBAC、企业身份和团队任务流；
  \item 企业数据连接器、持续数据更新和运营后台；
  \item 生产监控、告警、并发、成本治理和 SLA；
  \item 人工 gold 相关性集和更开放的质量评估；
  \item 多名目标运营用户的真实任务与业务效果证据。
\end{itemize}

主动说边界不会削弱项目，反而说明你知道原型、工程验证、用户验证和生产系统之间的区别。

\section{面试前自测}

\begin{enumerate}
  \item 能否不看文档画出一次完整请求的主路径和三条分支？
  \item 能否用自己的话解释为什么模型没有商品和工具执行权？
  \item 能否运行 synthetic 工作台并解释每个状态？
  \item 能否复现一个验证失败或证据不足分支？
  \item 能否说出 P5、P6 offline、P6 live 和 P7 各自能证明什么？
  \item 能否讲清一个真实 Bad Case 的 RED、修复和 GREEN？
  \item 能否在不夸大结果的情况下回答项目价值？
  \item 能否明确 P7 为什么仍为 \statusprogress？
\end{enumerate}

如果其中任何一项只能背诵，先回到对应章节做练习，再把它写进简历。

\section{证据指针}

本章的详细案例和问答依据 \filelink{docs/interview-casebook.md}，阶段状态以 \filelink{docs/master-plan.md} 和 \filelink{docs/current-stage-status.md} 为准。具体评测结论应回到 P5、P6 与 P7 的公开报告，不使用可能过期的口头数字替代同版本证据。
